\documentclass[11pt,a4paper,twoside]{tesis}
% SI NO PENSAS IMPRIMIRLO EN FORMATO LIBRO PODES USAR
%\documentclass[11pt,a4paper]{tesis}

\usepackage{graphicx}
\usepackage[utf8]{inputenc}
\usepackage[spanish]{babel}
\usepackage[left=3cm,right=3cm,bottom=3.5cm,top=3.5cm]{geometry}

\begin{document}

%%%% CARATULA

\def\autor{Autor}
\def\tituloTesis{Título del Trabajo de Especialización}
\def\runtitulo{Resumen}
\def\runtitle{Título del Trabajo}
%\def\director{Obi-Wan Kenobi}
%\def\codirector{Master Yoda}
\def\lugar{Buenos Aires, 2022}
\newcommand{\HRule}{\rule{\linewidth}{0.2mm}}
%
\thispagestyle{empty}

\begin{center}\leavevmode

\vspace{-2cm}

\begin{tabular}{l}
\includegraphics[width=3.6cm]{logouba.png}
\end{tabular}


{\large \sc Universidad de Buenos Aires\\Facultad de Ciencias Exactas y Naturales \\ Facultad de Ingeniería}

\vspace{6.0cm}

%\vspace{3.0cm}
%{
%\Large \color{red}
%\begin{tabular}{|p{2cm}cp{2cm}|}
%\hline
%& Pre-Final Version: \today &\\
%\hline
%\end{tabular}
%}
%\vspace{2.5cm}

\begin{huge}
\textbf{\tituloTesis}
\end{huge}

\vspace{2cm}

{\large Trabajo Final de Especialización en Explotación de Datos y Descubrimiento del Conocimiento}

\vspace{2cm}

{\Large \autor}

\end{center}

\vfill

{\large

%{Director: \director}

\vspace{.2cm}

%{Codirector: \codirector}

\vspace{.2cm}

\lugar
}

\newpage\thispagestyle{empty}


%%%% ABSTRACTS, AGRADECIMIENTOS Y DEDICATORIA
\frontmatter
\pagestyle{empty}
%\begin{center}
%\large \bf \runtitulo
%\end{center}
%\vspace{1cm}
\chapter*{\runtitulo}

\noindent Acá va un resumen del trabajo. Con el resumen se deberia poder tener una idea del trabajo en su totalidad. Desde los obetivos y los datos hasta los resultados y conclusiones. Se escibe al final. (aprox. 200 palabras).

\bigskip

\noindent\textbf{Palabras claves:} representativas del trabajo, los métodos, los datos (no menos de 5).

%\cleardoublepage
%%\begin{center}
%\large \bf \runtitle
%\end{center}
%\vspace{1cm}
\chapter*{\runtitle}

\noindent In a galaxy far, far away, a psychopathic emperor and his most trusted servant -- a former Jedi Knight known as Darth Vader -- are ruling a universe with fear. They have built a horrifying weapon known as the Death Star, a giant battle station capable of annihilating a world in less than a second. When the Death Star's master plans are captured by the fledgling Rebel Alliance, Vader starts a pursuit of the ship carrying them. A young dissident Senator, Leia Organa, is aboard the ship \& puts the plans into a maintenance robot named R2-D2. Although she is captured, the Death Star plans cannot be found, as R2 \& his companion, a tall robot named C-3PO, have escaped to the desert world of Tatooine below. Through a series of mishaps, the robots end up in the hands of a farm boy named Luke Skywalker, who lives with his Uncle Owen \& Aunt Beru. Owen \& Beru are viciously murdered by the Empire's stormtroopers who are trying to recover the plans, and Luke \& the robots meet with former Jedi Knight Obi-Wan Kenobi to try to return the plans to Leia Organa's home, Alderaan. After contracting a pilot named Han Solo \& his Wookiee companion Chewbacca, they escape an Imperial blockade. But when they reach Alderaan's coordinates, they find it destroyed - by the Death Star. They soon find themselves caught in a tractor beam \& pulled into the Death Star. Although they rescue Leia Organa from the Death Star after a series of narrow escapes, Kenobi becomes one with the Force after being killed by his former pupil - Darth Vader. They reach the Alliance's base on Yavin's fourth moon, but the Imperials are in hot pursuit with the Death Star, and plan to annihilate the Rebel base. The Rebels must quickly find a way to eliminate the Death Star before it destroys them as it did Alderaan (aprox. 200 palabras).

\bigskip

\noindent\textbf{Keywords:} War, Rebellion, Wookie, Jedi, The Force, Empire (no menos de 5). % OPCIONAL: comentar si no se quiere

%\cleardoublepage
%\chapter*{Agradecimientos}

\noindent  % OPCIONAL: comentar si no se quiere

%\cleardoublepage
%\hfill \textit{A mi persona favorita.}
  % OPCIONAL: comentar si no se quiere

%\cleardoublepage
\tableofcontents

\mainmatter
\pagestyle{headings}

%%%% ACA VA EL CONTENIDO DE LA TESIS

\chapter{Introducción}

\section{Descripción del problema y motivación ($>$ 1/2 carilla)}

Ejemplo de ecuación:
\begin{equation}
L = - \sum_{p \in P} \sum_{y \in Y} w_{py} \sum_{t \in T_{py}} y_t \log \hat y_t + (1-y_t) \log (1-\hat y_t)
\label{eq:loss}
\end{equation}

\section{Trabajos previos ($>$ 1 carilla)}
En la imagen \ref{fig:label} hay un ejemplo de como poner una figura:
\begin{figure}[h]
\centering
\includegraphics[width=0.5\columnwidth]{logouba.png}
\caption{Acá va una explicación de la figura. Recuerden usar colores solo si son explicativos, no olvidar nombrar correctamente los ejes y poner una explicación en el caption.}
\label{fig:label}
\end{figure}

\section{Objetivos ($>$ 1/2 carilla)}

\chapter{Materiales y Métodos}

\section{Datos ($>$ 1 carilla)} 
Descripción de los datos / obligatorio poner la fuente.

\section{Análisis exploratorio}

\chapter{Métodos ($>$ 3 carillas)}
Describir los métodos utilizados (a utilizar). 
Es bligatorio agregar un parrafo especificando a que materías se relacionan.

\section{Método I }
\section{Método II}
\section{Métricas}

\chapter{Experimentos (si los hay)}
Pueden ser comparaciones entre los métodos utilizados, por ejemplo. 

En la \ref{table:tab} se ve un ejemplo de tabla:
\begin{table}[h!]
\centering
\footnotesize
\begin{tabular}{lrrrrrr}
\hline
Phone &  Total & \% Errors &F1 &  EER \\
\hline
   EY &    441 &  13.83 & 0.80 & 0.12 \\
   JH &    178 &  39.89 & 0.81 & 0.18 \\
   AY &   1040 &   5.96 & 0.45 & 0.21 \\
    R &   1298 &  18.34 & 0.53 & 0.27 \\
\hline
\end{tabular}
\caption{Describir la tabla como para que alguien que agarra el trabajo no dependa del cuerpo del texto para entenderla}
\label{table:tab}
\end{table}

\chapter{Resultados ($>$ 3 carillas)}
Pueden ser preliminares. Es bligatorio tener algo hecho.


\chapter{Conclusiones ( $>$ 1 carilla)} 
Pueden ser preliminares. 
Si no hay mucho hecho se pueden discutir las dificultades a futuro.


%%%% BIBLIOGRAFIA
\backmatter
%\bibliography{tesis}

\end{document}
